% Chapter 1. Every number is a macro (generated/numbers.tex).
\chapter{Introduction}
\label{ch:intro}

\section{The silence after the user stops}

People taking turns in conversation leave short gaps between them
\cite{stivers2009universals}. A voice agent built as a cascade, with a speech
recognizer, a language model and a speech synthesizer, answers much later,
and later still when the whole cascade runs on a small device. On the device
studied here, a reactive pipeline answering spoken multi-step math problems
starts speaking a median of \TTFAMulti\,ms after the user stops.

The field's best-known remedy is to \emph{think while listening}: start the
language model before the user has finished, so that part of the reply is
ready when the endpoint arrives
\cite{predgen2025,ltsvoiceagent2026,epa2026,livemind2024}. Spoken language
models pursue the same idea inside one model
\cite{shanks2025,stitch2026,shih2025thinklisten}. Whatever the design, the
benefit depends on the compute spent while the user speaks being otherwise
idle.

\section{A device that pays for it}

On an edge system-on-chip with unified memory, the recognizer, the reasoner
and the synthesizer share one memory system and one set of cores. Nothing
waits for free. A speculative decode competes with whatever runs next to it,
and co-residence alone can slow a neighbour before any decode runs. Whether
thinking while listening helps on such a device is an empirical question, and
the answer depends on what the silence is made of.

On our device, the silence is mostly the recognizer's. Its final decode is the
largest stage after speech on every utterance length we measure:
\SttShareMulti\% of the stage sum on multi-step problems, against
\LlmShareMulti\% for the language model's first token. Speculating the reply
attacks a small stage. Speculating the \emph{transcript}, by committing what
the recognizer has already agreed on while the user is still speaking, attacks
the largest.

\section{Thesis statement}

\begin{quote}
On a unified-memory edge device the silence after the user stops belongs
mostly to the recognizer, not the language model. A rule that allocates the
listening window from measured costs spends it on the listener, and the gain
is measured.
\end{quote}

\section{Research questions}

\begin{description}
\item[RQ1] Where does the time between the end of speech and the first reply
  audio go on a unified-memory edge device? (Chapter~\ref{ch:cost})
\item[RQ2] What does speculation cost on shared silicon, for each consumer of
  the listening window? (Chapter~\ref{ch:cost})
\item[RQ3] What is reply speculation worth on a realistic spoken task, and
  could the pipeline tell when it would be? (Chapter~\ref{ch:value})
\item[RQ4] Can a per-turn rule allocate the window, what does it gain, and how
  close is it to a planner that reasons about the endpoint?
  (Chapter~\ref{ch:allocator})
\end{description}

\section{Contributions}

\begin{enumerate}
\item \textbf{A decomposition of the silence} on a unified-memory device, by
  stage and utterance length, showing that the recognizer's final decode is
  the largest stage at every length.
\item \textbf{The costs of speculation, priced on one axis.}
  \begin{itemize}
  \item An occupancy tax on the recognizer, through huge-page fallback.
  \item A reply draft still decoding at the endpoint costs \ThinkerOverlap{}.
  \item A partial hypothesis still decoding at the endpoint costs the final
    \ListenerC\,ms per second of overlap.
  \end{itemize}
\item \textbf{A precise negative on the thinker's side.} On spoken multi-step
  problems a reply draft becomes usable only near the end of the utterance,
  where it no longer fits before the endpoint. The registered test closed the
  value side.
\item \textbf{A rule that allocates the window, and its measured gain.} A
  feasibility gate and a freeze of the orphaned hypothesis let the listener cut
  the median paired time to first audio on multi-step problems by \PthreeRed{}
  ($n=\PthreeN$ items). The interval's lower end lies below the bar we
  registered. On short utterances the gate abstains at no detectable TTFA
  cost.
\item \textbf{A formalisation} of the allocation problem as a finite-horizon
  planner over the endpoint hazard. In simulation, calibrated from the logs,
  the gate is within \PlannerDP{} of the planner's optimum at the median.
\item \textbf{A methods record.} Every prediction was registered before its
  data and every correction is kept in place. Every number is produced by a
  script from the logs.
\end{enumerate}

\section{Method in brief}

All measurements were made on one Jetson Orin Nano Super, with recorded and
synthesized speech played into the pipeline in real time. Experiments were
pre-registered in a lab notebook before they ran, with their predictions,
bars and scoring rules. Where a pass was descriptive, or a scorer was written
after the data, the ledger says so. Corrections are appended, never edited in. A
claim registry ties every number in this thesis to the script and the runs
that produce it (Chapter~\ref{ch:methods}).

\section{Outline}

Chapter~\ref{ch:background} reviews cascaded voice agents, streaming
recognition, thinking while listening, and interference on shared hardware.
Chapter~\ref{ch:methods} describes the testbed and the methodology, including
the methods results. Chapter~\ref{ch:cost} measures the cost of thinking while
listening, and Chapter~\ref{ch:value} its value. Chapter~\ref{ch:allocator}
builds and measures the window allocator. Chapter~\ref{ch:discussion}
discusses the results and their limitations, and Chapter~\ref{ch:conclusion}
concludes. The appendices hold the capture-path hazards, the retractions log
and the full pre-registration ledger.
